\section{Introduction}
\label{sec:introduction}

Rapid damage assessment helps responders prioritize inspection and allocate scarce attention during disasters. Social-media reports often pair images with contemporaneous text, and CrisisMMD provides a widely used benchmark for this multimodal setting across seven 2017 disasters \citep{alam2018crisismmd}. Multimodal classifiers and newer VLM-based systems suggest that visual and linguistic signals can support disaster categorization, damage analysis, and report generation \citep{ofli2020multimodal,agarwal2020crisisdias,shetty2025disaster,chen2024disasteller}. Yet the same interface creates a trust boundary: accompanying text may clarify a scene, contradict it, or explicitly instruct a model to disregard what is visible.

Optimized image hijacks, typographic prompts, and multimodal attacks show that visual or externally supplied text can redirect VLM behavior \citep{bailey2024imagehijacks,cheng2024typographic,gong2025figstep,cao2025scenetap}. Cross-modal conflict studies further show that misleading text can override relevant visual evidence \citep{deng2025wordsvision,zhang2026misinformation}. What remains unclear is whether the same message has different safety consequences when delivered through the image, the accompanying text, or both, and whether malicious wording changes decisions beyond the instability caused by a benign overlay or prefix. Earlier crisis-domain adversarial work examined training for multimodal crisis classification \citep{chen2020crisisadv}, but it predates instruction-following VLMs and does not measure directional under-triage. We address this gap in an ordinal disaster-damage task, where moving a correct severe or mild judgment downward has a clear operational meaning.

We study this question through a controlled, paired robustness audit. Ground truth is the CrisisMMD image damage-severity annotation with three ordered labels: little/no, mild, and severe damage. The 720-source main cohort is class-balanced and globally separated from auxiliary cohorts by exact tweet/image identity and perceptual duplicate cluster. Each source is evaluated clean and under nine interventions: modality-matched benign controls plus direct or misleading malicious messages delivered through image, text, or both (Figure~\ref{fig:overlay}). Separate cohorts examine presentation style, relative size, text rhetoric, and nominal point size without changing the primary estimand. All experiments use the same fixed zero-shot rubric and deterministic decoding.

The design separates three questions: which delivery channel is most damaging, whether direct instructions and misleading claims behave differently, and how much of the change is attributable to malicious wording rather than to the intervention itself. Secondary cohorts vary presentation and text size, and the six-model panel shows whether the main pattern is confined to one system. Balanced-main accuracy spans 50.28--55.69\%, leaving 232--294 clean-correct mild/severe decisions per model. We therefore report clean competence separately from attack susceptibility. The primary rate counts initially correct mild/severe decisions that move downward, while retaining all 720 sources in the denominator; a conditional eligible-only rate and the corresponding upward rate provide complementary views.

Every malicious model--condition pair exceeds its modality-matched benign control. All 36 paired full-cohort risk differences are positive, their bootstrap intervals exclude zero, and every Holm-adjusted McNemar test is significant. The magnitude and most damaging channel nevertheless vary by model: image-only delivery dominates for some systems, joint delivery for others, and the two are similar elsewhere. Text-only effects are smaller but still exceed their benign counterparts. Upward shifts are rare (condition means at most 0.60\%), whereas severe-to-lower transitions can be large, locating the principal risk in under-triage rather than generic label instability.

Our main contribution is a paired evaluation design for directional harm in multimodal disaster assessment. It holds the source and assigned message constant across delivery channels, uses modality-matched benign controls to separate malicious influence from overlay or prefix instability, and reports both downward and upward movement in an ordered label space. Applied to duplicate-cluster-disjoint cohorts, the design reveals a common direction of failure across six models without assuming a common effect size or modality ranking. Presentation, relative-size, rhetoric, and point-size studies then test how far that conclusion extends beyond the main rendering and wording choices.
